\documentclass[aspectratio=169,10pt]{beamer}
\usetheme{default}
\setbeamertemplate{navigation symbols}{}
\usepackage{graphicx}\usepackage[T1]{fontenc}\usepackage{amsmath,amssymb}
\usepackage{listings}
\definecolor{ncsured}{HTML}{CC0000}\definecolor{ncsudk}{HTML}{990000}
\definecolor{ncsugray}{HTML}{6D6D6D}\definecolor{ncsuteal}{HTML}{008473}
\definecolor{codebg}{HTML}{F5F5F5}
\setbeamercolor{structure}{fg=ncsured}
\setbeamercolor{frametitle}{fg=white,bg=ncsured}
\setbeamercolor{title}{fg=ncsudk}
\setbeamercolor{itemize item}{fg=ncsured}
\setbeamercolor{itemize subitem}{fg=ncsured}
\setbeamerfont{frametitle}{series=\bfseries}
\setbeamertemplate{frametitle}{\nointerlineskip\vskip2pt%
  \begin{beamercolorbox}[wd=\paperwidth,ht=2.4ex,dp=1ex,leftskip=.3cm]{frametitle}%
  \small\bfseries\insertframetitle\end{beamercolorbox}}
\setbeamertemplate{footline}{\hbox{\begin{beamercolorbox}[wd=\paperwidth,ht=2.2ex,dp=1ex,leftskip=.3cm,rightskip=.3cm]{}%
  \tiny Optimal Market Making for Cryptocurrency Options \hfill Week 3 \hfill \insertframenumber\end{beamercolorbox}}}
\lstset{basicstyle=\ttfamily\scriptsize,backgroundcolor=\color{codebg},
  keywordstyle=\color{ncsured},commentstyle=\color{ncsugray},
  stringstyle=\color{ncsuteal},breaklines=true,columns=fullflexible,
  showstringspaces=false,frame=single,rulecolor=\color{ncsugray}}

\title{Week 3: The Avellaneda-Stoikov Model}
\subtitle{Optimal market making via stochastic control -- derived from scratch}
\date{North Carolina State University \,\textbar\, Summer 2026}
\titlegraphic{
\includegraphics[width=4.2cm]{/Users/dendisuhubdy/Github/ncsu/assets/ncsu_logo.png}}
\begin{document}
\begin{frame}[plain]\titlepage\begin{center}\tiny Industry Mentor: Dendi Suhubdy (Bitwyre)  |  Guest: Fenni Kang (Coincall)\end{center}\end{frame}
\begin{frame}{Learning Objectives}
\begin{itemize}
  \item Set up the Avellaneda-Stoikov control problem
  \item Follow the HJB derivation and understand every modeling choice
  \item Interpret the reservation price and the optimal half-spread
\end{itemize}
\end{frame}

\begin{frame}{Lecture 1: Setup \& Dynamics}
\begin{itemize}
  \item Mid-price: dS\_t = sigma dW\_t  (arithmetic Brownian motion)
  \item Fill intensity: lambda(delta) = A exp(-kappa delta)
  \item Wealth from fills; inventory q jumps by +/-1 on a fill
  \item Exponential (CARA) utility of terminal wealth
\end{itemize}
\end{frame}

\begin{frame}{Lecture 2: HJB, Ansatz, Closed Form}
\begin{itemize}
  \item Dynamic programming -\textgreater{} HJB partial differential equation
  \item Separation ansatz collapses HJB to ODEs in inventory
  \item Asymptotic (long-horizon) closed-form quotes
  \item Reservation price skews quotes against inventory
\end{itemize}
\end{frame}

\begin{frame}{Illustration}
\begin{center}\includegraphics[width=0.86\textwidth]{/Users/dendisuhubdy/Github/ncsu/weekly_presentations/figures/week03.png}\end{center}
\end{frame}

\section*{Derivation}
\begin{frame}{Derivation \& Rationale}
\[dS_t = \sigma\, dW_t\]

\small\textbf{\color{ncsured}Why:} Arithmetic (not geometric) Brownian motion is chosen because over a short market-making horizon price changes are small and roughly additive; it keeps the HJB linear in $S$ and makes the cash dynamics exact.
\end{frame}

\begin{frame}{Derivation \& Rationale}
\[\lambda(\delta) = A\, e^{-\kappa \delta}\]

\small\textbf{\color{ncsured}Why:} Fill intensity decays with quote distance $\delta$. The exponential form is the simplest law consistent with 'further from mid $\Rightarrow$ fewer fills'; $A$ sets the base rate, $\kappa$ the decay. It also yields a closed-form optimum.
\end{frame}

\begin{frame}{Derivation \& Rationale}
\[\max\ \mathbb{E}\!\left[-e^{-\gamma\,(X_T + q_T S_T)}\right]\]

\small\textbf{\color{ncsured}Why:} CARA (exponential) utility is used because constant absolute risk aversion makes the value function factor out wealth, so the optimal quotes do not depend on the cash level -- exactly what we want for a market maker.
\end{frame}

\begin{frame}{Derivation \& Rationale}
\[\partial_t V + \tfrac12\sigma^2 \partial_{SS} V + \max_{\delta^b}\lambda(\delta^b)\big[V(q{+}1,x{-}S{+}\delta^b)-V\big]+ \max_{\delta^a}\lambda(\delta^a)\big[V(q{-}1,x{+}S{+}\delta^a)-V\big]=0\]

\small\textbf{\color{ncsured}Why:} The HJB equation: the diffusion term prices inventory risk; each $\max$ term is the dealer optimally choosing the bid/ask offset, trading fill probability against captured spread.
\end{frame}

\begin{frame}{Derivation \& Rationale}
\[V(t,x,S,q) = -e^{-\gamma(x+qS)}\,\theta(t,q)\]

\small\textbf{\color{ncsured}Why:} This ansatz is not a guess from nowhere: under CARA $+$ arithmetic dynamics the value function inherits a multiplicative structure across $(x, qS)$, so all the wealth/price dependence is exponential and only $\theta(t,q)$ remains to solve.
\end{frame}

\begin{frame}{Derivation \& Rationale}
\[r(t,q) = S - q\,\gamma\,\sigma^2 (T-t)\]

\small\textbf{\color{ncsured}Why:} The reservation (indifference) price: it sits BELOW mid when you are long ($q>0$) so you quote to sell down inventory. The skew scales with risk aversion $\gamma$, variance $\sigma^2$, and time left -- more risk $\Rightarrow$ more skew.
\end{frame}

\begin{frame}{Derivation \& Rationale}
\[\delta^* = \tfrac{1}{\gamma}\ln\!\big(1+\tfrac{\gamma}{\kappa}\big) + \tfrac{\gamma\sigma^2}{2}(T-t)\]

\small\textbf{\color{ncsured}Why:} The optimal half-spread = an order-flow term $\tfrac1\gamma\ln(1+\gamma/\kappa)$ (how hard fills are to get) plus an inventory-risk term $\tfrac{\gamma\sigma^2}{2}(T-t)$. You widen when risk-averse, when volatile, and when far from the horizon.
\end{frame}

\begin{frame}{Getting the Data from Coincall}
\begin{itemize}
  \item This week uses simulated dynamics (no live data needed).
  \item For realistic sigma, estimate it from Coincall perp klines:
  \begin{itemize}\item \texttt{\small GET /open/option/market/kline/history/v1/\{name\}?period=m1  (or futures klines)}\end{itemize}
  \item You will calibrate A, kappa from real fills in Week 5.
\end{itemize}
\vspace{2pt}\tiny\textit{Cross-check exact paths at https://docs.coincall.com/}
\end{frame}

\begin{frame}[fragile]{Code --- week3\_as\_quotes.py}
\begin{lstlisting}
import torch

def as_quotes(S, q, gamma, sigma, kappa, tau):
    """Asymptotic Avellaneda-Stoikov bid/ask. Vectorized over inventory q."""
    reservation = S - q * gamma * sigma**2 * tau
    half_spread = (1.0/gamma) * torch.log1p(gamma/kappa) \
                  + 0.5 * gamma * sigma**2 * tau
    return reservation - half_spread, reservation + half_spread  # bid, ask

S = torch.tensor(100.0); q = torch.arange(-10, 11).float()
bid, ask = as_quotes(S, q, gamma=0.1, sigma=2.0, kappa=1.5, tau=1.0)
print("inventory ->", q.tolist())
print("bid       ->", bid.round(decimals=2).tolist())
\end{lstlisting}
\end{frame}
\begin{frame}{Glossary --- Key Terms}
\footnotesize\begin{description}
\item[\color{ncsured}Reservation price] inventory-adjusted indifference price that quotes are centered on
\item[\color{ncsured}Half-spread] distance from the reservation price to each quote
\item[\color{ncsured}HJB equation] the PDE characterizing the optimal value function
\item[\color{ncsured}Ansatz] an assumed functional form that reduces the HJB to ODEs
\item[\color{ncsured}CARA utility] constant absolute risk aversion, U(w) = -exp(-gamma w)
\item[\color{ncsured}Intensity lambda(delta)] Poisson fill rate as a function of quote distance delta
\item[\color{ncsured}Risk aversion gamma] parameter setting inventory skew and spread width
\end{description}
\end{frame}

\begin{frame}{Reading}
\begin{itemize}
  \item Avellaneda-Stoikov (2008); Cartea-Jaikumar-Penalva, Ch. 10
\end{itemize}
\end{frame}

\end{document}